\documentclass[11pt]{article}
\usepackage[a4paper,margin=1in]{geometry}
\usepackage{amsmath,amssymb,mathtools,bm}
\usepackage{enumitem,booktabs,array}
\usepackage{tcolorbox}
\usepackage{hyperref}
\usepackage[protrusion=true,expansion=false]{microtype}
\hypersetup{colorlinks=true,linkcolor=blue,urlcolor=blue,citecolor=blue}
\setlist[itemize]{topsep=2pt,itemsep=2pt}
\setlist[enumerate]{topsep=2pt,itemsep=2pt}
\newtcolorbox{keybox}{colback=blue!3!white,colframe=blue!40!black,boxrule=0.4pt,arc=1mm}
\newtcolorbox{alertbox}{colback=red!3!white,colframe=red!40!black,boxrule=0.4pt,arc=1mm}
\newtcolorbox{answerbox}{colback=green!3!white,colframe=green!40!black,boxrule=0.4pt,arc=1mm}
\newtcolorbox{drillbox}{colback=yellow!5!white,colframe=orange!60!black,boxrule=0.4pt,arc=1mm}
\newcommand{\EV}{\mathbb{E}}
\newcommand{\Var}{\mathrm{Var}}
\newcommand{\Cov}{\mathrm{Cov}}
\newcommand{\blank}[1]{\underline{\hspace{#1}}}

\title{W04-03: Mian \& Sufi (2009, QJE) \\
\large ``The Consequences of Mortgage Credit Expansion: Evidence from the U.S. Mortgage Default Crisis'' --- Active Recall Workbook}
\author{Banking \& Risk Management --- Exam Prep}
\date{\today}
\begin{document}\maketitle

\begin{keybox}
\textbf{Paper in one sentence.} ZIP codes with the \emph{largest expansion of mortgage credit to ex-ante high-denial-rate borrowers} between 2002--2005 experienced the largest subsequent default rates 2006--2008, despite the \emph{decline} in borrower-income growth in those same ZIPs --- direct evidence that a credit-supply outward shift (not a demand/income story) drove the subprime boom and the default crisis.

\textbf{Placement.} The companion to Keys et al.\ (W04-02): KMSV gives the micro screening-vs-securitisation channel; MS shows the macro geography of the credit supply expansion and its default consequences. Keystone of the household-debt literature.
\end{keybox}

\section*{Round 1: Complete Walk-Through}

\subsection*{1.1 Empirical puzzle}
Between 2002 and 2005, mortgage credit expanded most in ZIP codes that had \emph{high denial rates} in the 1990s (subprime-prone neighbourhoods). This is striking because ZIP-level income growth in those neighbourhoods was \emph{negative or zero} --- the outward shift in credit cannot be demand-driven.

\subsection*{1.2 Data}
HMDA (Home Mortgage Disclosure Act) ZIP-code panel: originations, denials, borrower-income measures. Default and price data from LoanPerformance and Zillow. Sample: $\approx 3{,}000$ ZIP codes across 2001--2008.

\subsection*{1.3 Identification --- supply vs.\ demand}
The MS insight: if the credit expansion were demand-driven (more borrower income or wealth), origination and income should move \emph{together}. MS document the opposite in subprime ZIPs: originations up, income flat or down. This \emph{negative correlation} of origination growth and income growth is the signature of a supply shock.
\[
\Delta\ln\text{Orig}_{z,02\text{-}05}=\alpha + \beta\cdot\text{HighDenial90s}_z + \gamma\cdot\Delta\ln\text{Income}_{z,02\text{-}05}+\varepsilon_z.
\]
The coefficient $\beta$ is positive, and the simple correlation with income growth is \emph{negative} in subprime ZIPs.

\subsection*{1.4 Default consequences}
Same-ZIP regressions in 2006--2008:
\[
\text{DefaultRate}_{z}= a + b\cdot\Delta\ln\text{Orig}_{z,02\text{-}05}+\text{controls}+u_z.
\]
High-credit-growth subprime ZIPs have by far the highest default rates in the subsequent bust; the relationship is linear and strong.

\subsection*{1.5 Mechanism}
\begin{itemize}
\item Securitisation finance (via private-label MBS) flows disproportionately to subprime ZIPs from 2002.
\item MS show that non-GSE securitisation is the marginal funding source.
\item House prices rise in these ZIPs, then crash more dramatically (housing-bubble amplification).
\end{itemize}

\subsection*{1.6 Caveats}
\begin{alertbox}
\begin{itemize}
\item ZIP-level analysis may mask within-ZIP heterogeneity; MS supplement with loan-level HMDA.
\item The counterfactual ``what would defaults look like without supply shift'' is not observed; MS use cross-sectional variation.
\item Reverse causality --- defaults driving credit policy --- is unlikely given the timing (2002--2005 expansion precedes 2006--2008 defaults).
\end{itemize}
\end{alertbox}

\section*{Round 2: Level 1 Fill-in}
\begin{enumerate}
\item Expansion period: \blank{1cm}--\blank{1cm}; default period: \blank{1cm}--\blank{1cm}.
\item Treatment: ZIPs with high ex-ante \blank{1.5cm} rate.
\item Demand-side test: origination growth vs.\ \blank{1.5cm} growth is \blank{1cm}.
\item Marginal funding source: non-GSE \blank{1.5cm}.
\item Outcome: ZIP-level default rate 2006--\blank{1cm}.
\end{enumerate}

\section*{Round 3: Level 2 Prompts}
\begin{enumerate}
\item Why does the sign of origination-vs-income correlation identify supply vs.\ demand?
\item Relate MS to Keys et al.\ (2010) and Gorton--Metrick (2012). How does the story fit together?
\item How would MS's evidence change if non-GSE securitisation had been absent?
\item Discuss implications for macroprudential policy.
\end{enumerate}

\section*{Round 4: Minimal Scaffolding}
From a blank page: (i) empirical puzzle, (ii) data and geography, (iii) supply vs.\ demand argument, (iv) default consequences, (v) two caveats.

\section*{Round 5: Blank Page}
Essay: the subprime boom as a credit-supply shift --- MS evidence and its implications.

\vspace{6pt}
\begin{drillbox}
\textbf{Speed Drill.} (1) Expansion years. (2) Treatment variable. (3) Sign of origination--income correlation in subprime ZIPs. (4) Default years. (5) Funding source identified.
\end{drillbox}

\section*{Quick Reference \& Answer Key}
\begin{answerbox}
\begin{itemize}
\item \textbf{Puzzle.} Subprime ZIPs: mortgage credit $\uparrow$ 2002--05 while income $\downarrow$.
\item \textbf{Inference.} Supply shift, not demand, drives the expansion.
\item \textbf{Default.} Same ZIPs have the highest 2006--08 default rates.
\item \textbf{Channel.} Non-GSE (private-label) securitisation as the marginal funder.
\end{itemize}
\end{answerbox}

\begin{keybox}
\textbf{30-second exam pitch.} Mian \& Sufi (2009) pin down the supply-side origins of the US mortgage default crisis. Between 2002 and 2005, mortgage credit expanded most in ZIP codes that had been the \emph{hardest} to borrow in during the 1990s --- high ex-ante denial-rate neighbourhoods. At the same time, income growth in those ZIPs was flat or negative, ruling out a demand-driven expansion. The marginal funding source was non-GSE (private-label) securitisation. These same ZIPs went on to have the highest default rates in 2006--2008. The paper is the macro counterpart to Keys et al.\ (2010): it shows that the screening-lax securitisation channel manifested in a geographic pattern where credit expansion was decoupled from income fundamentals, and the reversal drove the household-debt collapse.
\end{keybox}

\end{document}
